% Source: Independent optimization checklist.
\subsection{DP 优化的检查顺序}
\label{dp:optimization-basics}
\subsubsection*{原理}

先保留明确的直接递推，再确定哪一部分耗时。常见方向如下：
\begin{itemize}
 \item 前缀和或前缀最值：把一段状态的重复求和、求最值预处理掉。
 \item 滚动数组：减少存储层数，核对更新顺序，不自动降低转移时间。
 \item 单调队列：候选区间滑动，比较函数在进入队列时就能确定。
 \item 单调栈：按当前贡献把候选分段，相邻段的合并只发生一次。
 \item 斜率／凸包：将候选写成直线，先检查斜率与查询横坐标的单调性。
 \item 分治：最优决策位置随状态单调，利用它缩小枚举范围。
\end{itemize}
每种优化都依赖额外性质。只看到“枚举 j 取最小”不能自动使用凸包或分治；应先推导候选函数、有效范围、何时插入、是否过期，以及决策是否真的单调。验证新优化时，与小规模直接 DP 对照即可，无需每次重跑所有无关算法。
